\documentclass[10pt,a4paper]{amsart}
\usepackage[margin=19mm]{geometry}
\usepackage{amsmath,amssymb,mathtools}
\usepackage[scr=rsfs,cal=euler]{mathalfa}
\usepackage{xfrac}
\usepackage[unicode,hidelinks,bookmarks=false]{hyperref}
\newcommand{\ie}{i.e.\ }
\newcommand{\cf}{cf.\ }
\newcommand{\resp}{respectively\ }
\newcommand{\Subsection}{\subsection}
\providecommand{\nobreakdash}{\nobreak\hbox{-}}
\setlength{\emergencystretch}{2em}
\multlinegap0pt
\usepackage{bm}
\usepackage{bbm}
\usepackage{dsfont}
\usepackage{xspace}
\usepackage{cancel}
\usepackage{mathtools}
\usepackage{amsthm}
\makeatletter
\@ifundefined{lemma}{\newtheorem{lemma}{Lemma}}{}
\@ifundefined{proposition}{\newtheorem{proposition}{Proposition}}{}
\makeatother
\title[Chain controllability of linear control systems]{Chain controllability of linear control systems}
\author{Fritz Colonius, Alexandre J. Santana, Eduardo C. Viscovini}
\date{Source: arXiv:2312.14595v1 (CC BY 4.0)}
\begin{document}
\maketitle

\noindent\emph{Source and licence.} Chain controllability of linear control systems. Version arXiv:2312.14595v1 (CC BY 4.0), distributed under the Creative Commons Attribution 4.0 International licence, as recorded in the arXiv licence metadata for this exact version and shown on the abstract page https://arxiv.org/abs/2312.14595v1. Full rights notice: licenses/arxiv-2312.14595-CC-BY-4.0.txt.\par\smallskip

\noindent\textit{Editorial scope.} Counted unit: the Lemma whose source label is \texttt{Lemma2\_Eduardo} in the article (source0.tex lines 922--926, source label \texttt{Lemma2\_Eduardo}), delivered together with the article's own complete original proof of it (source0.tex lines 928--1035, one \texttt{proof} environment of 108 lines). No mathematics is written, repaired or re-derived: statement and proof are the article's own text line for line. Required context retained uncounted: LemCLSAtReverseTime (lines 219--223). Every definition, display and label of those lines is retained unchanged. Omitted from this package and not redistributed: 1--218, 224--921, 927--927, 1036--2224. Nothing inside a retained excerpt is omitted or altered: every retained body is delivered exactly as the article's lines are written, so no omission receipt of any kind accompanies this package. Citations: the retained excerpts carry no citation command at all, so no bibliography entry of the article is transcribed and no reference is redistributed. Reference adaptation: none was needed and none was made; every reference inside the delivered unit resolves inside it, so no unresolved reference or citation is produced. Printed slips kept literal: no printed word, symbol, hypothesis or number of the counted statement or of its proof was altered. Layout and class: the article's own class is \texttt{siamltex} with packages including amsfonts, amsmath, amssymb, graphicx, enumerate, color, hyperref; the wrapper is the lane's pinned \texttt{amsart} editorial wrapper at 10pt on A4, which restates the theorem-like environments the retained text needs and adds the article's own macro definitions that the retained text needs (none), with extra wrapper packages (bm, bbm, dsfont, xspace, cancel, mathtools). The delivered numbering is the unit's own. Assets: no retained span contains a figure environment, a graphics-inclusion command, a PDF-page inclusion, a table or a TikZ picture; no image file is copied into this package, the package declares no asset dependency and no image file is redistributed with it. Licence: version arXiv:2312.14595v1 of this article is distributed under the Creative Commons Attribution 4.0 International licence, as recorded in the arXiv licence metadata for this exact version and shown on the abstract page https://arxiv.org/abs/2312.14595v1. A full rights notice is delivered at licenses/arxiv-2312.14595-CC-BY-4.0.txt. Characters: every retained excerpt is pure ASCII, so the delivered engine, XeLaTeX with its default Latin Modern text font, reports no missing character.
\begin{proposition}
\label{LemCLSAtReverseTime}For any flow $\Psi$ on $X$ the forward and backward
chain limit sets are related by $y\in\Omega(x)$ if and only if $x\in
\Omega^{\ast}(y)$.
\end{proposition}

\begin{lemma}
\label{Lemma2_Eduardo}If all eigenvalues of the matrix $A$ have null real
parts, then for all $x,y\in\mathbb{K}^{n}$ and all $\varepsilon,T>0$ there is
an $(\varepsilon,T)$-chain of the flow $\psi$ from $x$ to $y$.
\end{lemma}

\begin{proof}
We first prove this lemma for the complex case $\mathbb{K}=\mathbb{C}$ using
induction over the dimension $n$. The lemma holds trivially for $n=0$ since
$\mathbb{C}^{0}$ consists of a single point. Assume that the assertion holds
for all $A\in\mathbb{C}^{(n-1)\times(n-1)}$. Since $\mathbb{C}$ is
algebraically closed, $A$ has some eigenvector $v$ associated to an eigenvalue
$\mu\in\mathbb{C}$ and, by hypothesis, $\operatorname{Re}\mu=0$. There is
$r>0$ such that $r\mu=2j\pi\imath$ for some $j\in\mathbb{Z}$. It follows that
for all $\lambda\in\mathbb{C}$%
\[
e^{rA}\lambda v=\lambda e^{2j\pi\imath}v=\lambda v.
\]
Thus the restriction of $\psi$ to the linear span $\left\langle v\right\rangle
$ of $v$ is periodic. We claim that for all $\lambda_{1}v,\lambda_{2}%
v\in\left\langle v\right\rangle ,\lambda_{1},\lambda_{2}\in\mathbb{C}$, and
all $\varepsilon,T>0$ there is an $(\varepsilon,T)$-chain from $\lambda_{1}v$
to $\lambda_{2}v$: In fact, consider periodic solutions through $\lambda
_{1}v+\frac{i}{j}(\lambda_{2}-\lambda_{1})v,i\in\{0,1,\ldots,j\}$, where
$j\in\mathbb{N}$ is large enough such that%
\[
\left\Vert \lambda_{1}v+\frac{i}{j}(\lambda_{2}-\lambda_{1})v-\left[
\lambda_{1}v+\frac{i+1}{j}(\lambda_{2}-\lambda_{1})v\right]  \right\Vert
=\frac{1}{j}\left\vert \lambda_{2}-\lambda_{1}\right\vert \left\Vert
v\right\Vert <\varepsilon.
\]
Staying on each periodic solution long enough such that the time is greater
than $T$ and jumping from one periodic solution to the next yields an
$(\varepsilon,T)$-chain from $\lambda_{1}v$ to $\lambda_{2}v$.

Now let $W\subset\mathbb{C}^{n}$ be a subspace such that $\mathbb{C}%
^{n}=W\oplus\left\langle v\right\rangle $. There are linear transformations
$A_{1}:W\rightarrow W$ and $A_{2}:W\rightarrow\left\langle v\right\rangle $
such that $Aw=A_{1}w+A_{2}w$. We obtain for $w\in W$,%
\[
\frac{d}{dt}\psi(t,w)=A\psi(t,w)=A_{1}\psi(t,w)+A_{2}\psi(t,w)\text{ with
}A_{2}\psi(t,w)\in\langle v\rangle.
\]
The variation-of-parameters formula for $A_{1}$ shows that%
\[
\psi(t,w)=e^{At}w=e^{tA_{1}}w+\widehat{\psi}(t,w)\text{ with }\widehat{\psi
}(t,w):=\int_{0}^{t}e^{(t-s)A_{1}}A_{2}\psi(s,w)ds.
\]
On the other hand, $\psi(t,\lambda v)=e^{At}\lambda v=e^{t\mu}\lambda v$ and
hence by linearity%
\begin{equation}
\psi(t,w+\lambda v)=e^{tA_{1}}w+\widehat{\psi}(t,w)+e^{t\mu}\lambda v.
\label{e_1}%
\end{equation}


Let $x\in\mathbb{C}^{n}$. We will show that $0\in\Omega(x)$ and $x\in
\Omega(0)$. This implies the assertion, since any two points $x,y$ can be
connected by concatenating a chain from $x$ to $0$ with a chain from $0$ to
$y$.

Let $\varepsilon,T>0$ and write $x=w\oplus\lambda v\in\mathbb{C}^{n}$. By the
induction hypothesis for all $\varepsilon,T>0$ there is an $(\varepsilon
,T)$-chain from $w$ to $0$ given by $w_{0}=w,w_{1},\ldots,\allowbreak w_{k}=0$
and $T_{0},T_{1},\ldots,T_{k-1}\geq T$ with%
\[
\left\Vert e^{T_{i}A_{1}}w_{i}-w_{i+1}\right\Vert <\varepsilon\text{ for
}i=0,\ldots,k-1.
\]
Next we construct an $(\varepsilon,T)$-chain from $x=w\oplus\lambda
v\in\mathbb{C}^{n}$ to $0$. Define recursively%
\[
v_{0}=\lambda v,v_{i+1}=\widehat{\psi}(T_{i},w_{i})+e^{T_{i}\mu}v_{i}%
\in\left\langle v\right\rangle \text{ for }i=1,\ldots,k-1.
\]
Hence by (\ref{e_1})%
\begin{align}
&  e^{T_{i}A}(w_{i}+v_{i})-(w_{i+1}+v_{i+1})\label{e_2}\\
&  =e^{T_{i}A_{1}}w_{i}+\widehat{\psi}(T_{i},w_{i})+e^{T_{i}\mu}v_{i}%
-w_{i+1}-\widehat{\psi}(T_{i},w_{i})-e^{T_{i}\mu}v_{i}=e^{T_{i}A_{1}}%
w_{i}-w_{i+1}.\nonumber
\end{align}
For $x_{i}:=w_{i}+v_{i},i=0,\ldots,k$, it follows that%
\[
\left\Vert e^{T_{i}A}x_{i}-x_{i+1}\right\Vert =\left\Vert e^{T_{i}A_{1}}%
w_{i}-w_{i+1}\right\Vert <\varepsilon.
\]
Thus $x_{0}=x,x_{1},\ldots,x_{k}=v_{k}$ and $T_{0},T_{1},\ldots,T_{k-1}\geq T$
is an $(\varepsilon,T)$-chain from $x$ to $v_{k}\in\left\langle v\right\rangle
$. Concatenating this with an $(\varepsilon,T)$-chain from $v_{k}$ to $0$ in
$\left\langle v\right\rangle $ we obtain the desired chain.

The inclusion $x\in\Omega(0)$ follows from considering the time reversed
system and Proposition \ref{LemCLSAtReverseTime}.

Now assume $\mathbb{K}=\mathbb{R}$. The previous result yields that for all
$x,y\in\mathbb{R}^{n}\subset\mathbb{C}^{n}$ and for all $\varepsilon,T>0$
there are $(\varepsilon,T)$-chains in $\mathbb{C}^{n}$ from $x$ to $y$.
Consider such a chain given by $z_{0}=x,z_{1},\ldots,z_{k}=y$ in
$\mathbb{C}^{n}$ and $T_{0},\ldots,T_{k-1}\geq T$ with%
\[
\left\Vert e^{T_{i}A}z_{i}-z_{i+1}\right\Vert <\varepsilon\text{ for
}i=0,\ldots,k-1.
\]
Since the entries of $A$ are real, it follows that for all $i$%
\[
\left\Vert e^{T_{i}A}\operatorname{Re}z_{i}-\operatorname{Re}z_{i+1}%
\right\Vert =\left\Vert \operatorname{Re}\left(  e^{T_{i}A}z_{i}%
-z_{i+1}\right)  \right\Vert \leq\left\Vert e^{T_{i}A}z_{i}-z_{i+1}\right\Vert
<\varepsilon.
\]
This shows that one obtains an $(\varepsilon,T)$-chain on $\mathbb{R}^{n}$
from $z_{0}=x$ to $z_{k}=y$.
\end{proof}

\end{document}
